\section{Conclusion} \label{sec:conclusion}

This thesis investigated the users' mental models of \ac{LLM}-based systems in collaborative problem-solving scenarios, and the effect of explanations on the mental states users attribute to them. Using a grounded theory approach \parencite{Gioia2013}, the research derived \acp{DP} for explainable \ac{LLM}-based systems that support accurate mental models from a qualitative study with 12 participants.

The data analysed in this thesis were collected during a user study in which participants solved maths problems with the assistance of an explainable \ac{LLM}-based system. The study was a mixed-methods experiment, comparing a group with explanations to a group without explanations. The think-aloud protocols were analysed in two cycles of coding \parencite{Saldana2015}: (1) process coding to foreground the participants' actions with the system, and (2) pattern coding to reveal common behavioural patterns across participants. To retain grounding in the data, in vivo codes were used alongsidethe process and pattern codes.

The codes were then assembled into themes and aggregate dimensions. The themes and dimensions revealed that the participants' interactions with the system followed a cyclical pattern of decision to engage, information exchange, and appraisal of the system's outputs. Appraisal revealed that participants attributed a high level of competence to the system compared to themselves. Attributions of mental states to the system were spread across this circle and mostly enacted rather than explicitly voiced. Participants built mental models around assumed system knowledge and capabilities, and revised them when the system's behaviour conflicted with their expectations.

Using the meta-requirements and design requirements formulated based on a literature review \parencite{Boell2014} centred in the field of \ac{HAII}, \acp{DP} were derived from the themes and dimensions. These initial \acp{DP} were evaluated through expert interviews with \ac{HCI} researchers and an industry practitioner. Expert feedback indicated that the \acp{DP} were largely effective and applicable, but needed revision to address comprehension challenges. Experts repeatedly highlighted sycophantic behaviour as a key challenge for \ac{LLM}-based systems, suggesting that the \acp{DP} could help to mitigate these issues. The expert feedback was incorporated into the final version of the \acp{DP}, resulting in one \ac{DP} being split due to overload. The final \acp{DP} were presented in tabular form \parencite[cf.][]{Gregor2020}.

While the design theory as a whole was evaluated positively, a portion of the design space remained uncovered due to a lack of empirical data. The data analysed in the study did not allow inferences about the effect of explanations on mental models. The resulting gap was filled with \acp{DP} derived from the literature review. Future research should focus on the evaluating design theory and the additional theory-based \acp{DP}.
